\section{非微扰重正化与匹配}\label{sec:matching}

为恢复准 PDF 矩阵元素的连续极限，需要在格点上做非微扰重正化以处理线性与对数的紫外发散。在本工作中，我们采用文献~\cite{Stewart:2017tvs,Chen:2017mzz}中详述的 RI/MOM 方案，并用一圈匹配系数~\cite{Stewart:2017tvs}把结果匹配到 $\overline{\rm MS}$ PDF。

\subsection{格点上的 RI/MOM 重正化}\label{subsec:rimom}

空间关联函数 $O_\Gamma(z)$ 在连续理论的坐标空间中已被证明是可乘重正化的~\cite{Ji:2017oey,Ishikawa:2017faj}，这使得可以在 RI/MOM 方案~\cite{Martinelli:1994ty}中进行重正化。

对每个 $z$ 值，RI/MOM 重正化因子 $Z$ 通过要求准 PDF 算符矩阵元素在给定动量的非壳夸克态上的圈图修正为零来确定：
\begin{align}\label{eq:Z}
Z(z,p^R_z,a^{-1},\mu_R)=\left.\frac{\sum_s\langle p,s|O_{\gamma^t}(z)|p,s\rangle}{\sum_s\langle p,s|O_{\gamma^t}(z)|p,s\rangle_{\rm tree}}\right|_{\tiny\begin{matrix}p^2=-\mu_R^2 \\ \!\!\!\!p_z=p^R_z\end{matrix}}\,.
\end{align}
裸矩阵元素 $\sum_s\langle p,s|O_{\gamma^t}(z)|p,s\rangle$ 将在格点上由 $O_{\gamma^t}(z)$ 的砍腿格林函数 $\Lambda_{\gamma^t}(p,z)$ 结合狄拉克矩阵的投影算符 ${\cal P}$ 计算：
\begin{equation}
\sum_s\langle p,s|O_{\gamma^t}(z)|p,s\rangle = \mbox{Tr}\left[ \Lambda_{\gamma^t}(z,p){\cal P}\right]\,.
\end{equation}
由于 $O_\Gamma(z)$ 中洛伦兹协方差的破坏，RI/MOM 减除依赖于 $\mu_R$ 和 $p_z^R$ 两个标度。因此重正化因子 $Z(z,p^R_z,a^{-1},\mu_R)$ 既依赖于格距，也依赖于两个 RI/MOM 标度 $\mu_R$ 与 $p^R_z$。

基于 $O_\Gamma(z)$ 在格点上的对称性，砍腿格林函数 $\Lambda_{\gamma^t}(p,z)$ 不仅正比于树图结果 $\gamma^t$，还包含另外两个独立的洛伦兹结构：
\begin{equation}\label{eq:ME_decomposition}
\Lambda_{\gamma^t}(p,z)=\widetilde{F}_t(p,z) \gamma^t  + \widetilde{F}_{z}(p,z)\frac{p_t\gamma^z}{p_z} +\widetilde{F}_p(p,z)\frac{p_t\slashed{p}}{p^2}\,.
\end{equation}
上式中的 $\widetilde{F}_i$ 是在超立方群 $H(4)$ 下不变的独立形状因子。根据式~(\ref{eq:ME_decomposition})，RI/MOM 重正化因子 $Z$ 也将依赖于投影算符 ${\cal P}$。可以选择只挑出 $\widetilde{F}_t$~\cite{Stewart:2017tvs}，我们称之为最小投影（minimal projection）。该投影形式最简单，却捕获了 $\Lambda_{\gamma^t}(p,z)$ 中的全部紫外发散。另一种可选方案是取 ${\cal P} = \slashed p/(4p^t)$，我们称之为 $\slashed p$ 投影。采用最小投影与 $\slashed p$ 投影的重正化因子 $Z$ 定义为
\begin{align}\label{eq:def_renorm}
&Z_{mp}(z, p^R_z, a^{-1}, \mu_R)\equiv \widetilde{F}_t(p,z)\Big|_{\tiny\begin{matrix}p^2=-\mu_R^2 \\ \!\!\!\!p_z=p^R_z\end{matrix}},\\
&Z_{\slashed{p}}(z, p^R_z, a^{-1}, \mu_R)\nonumber\\
&\equiv \Big[\widetilde{F}_t(p,z)+\widetilde{F}_z(p,z)+\widetilde{F}_p(p,z)\Big]\bigg|_{\tiny\begin{matrix}p^2=-\mu_R^2 \\ \!\!\!\!p_z=p^R_z\end{matrix}}.
\end{align}

由格点计算得到的坐标空间裸核子矩阵元素
\begin{align}
\widetilde{h}(z,P_z,a^{-1}) = {1\over 2P^0} \langle P |O_{\gamma^t}(z)| P \rangle \nn
\end{align}
按如下方式重正化：
\begin{align} \label{eq:rimomh}
&\widetilde{h}_R(z,P_z, p^R_z,\mu_R)\nn\\
&=\left.Z^{-1}(z,p^R_z,a^{-1},\mu_{\tiny R})\widetilde{h}(z,P_z,a^{-1})\right|_{a\to 0} .
\end{align}
这里 $\widetilde{h}_R(z,P_z, p^R_z,\mu_R)$ 是重正化矩阵元素的连续极限。于是，RI/MOM 方案中的准 PDF $\widetilde{q}_R(x,P_z,p^R_z,\mu_R)$ 由 $\widetilde{h}_R(z,P_z, p^R_z,\mu_{\tiny R})$ 的傅里叶变换得到：
\begin{align}
\widetilde{q}_R(x,P_z, p^R_z,\mu_R) = P_z\int {dz\over 2\pi}\ e^{ix P_zz}\widetilde{h}_R(z,P_z, p^R_z,\mu_R).\nn\\
\end{align}
在 RI/MOM 方案中，$\widetilde{h}_R(z,P_z, p^R_z,\mu_R)$ 与 $\widetilde{q}_R(x,P_z, p^R_z,\mu_R)$ 不再依赖紫外调节参数，从而准 PDF 与 $\overline{\rm MS}$ PDF 之间的一步匹配可在维数正规化的连续理论中进行~\cite{Stewart:2017tvs}。

准 PDF 最终都要匹配到同一个 $\overline{\text{MS}}$ PDF，因此采用 $Z_{mp}$ 与 $Z_{\slashed p}$ 的两种投影应当给出相同结果。然而匹配系数只能在固定圈阶计算，故对投影算符的残余依赖不可避免。

用同样的逻辑可以得出结论：准 PDF 对 RI/MOM 标度 $\mu_R$ 与 $p_z^R$ 的依赖也应被匹配系数完全抵消。任何固定阶的匹配计算都不可避免地导致最终 PDF 结果存在残余的 $\mu_R$、$p_z^R$ 与 $P_z$ 依赖。这些依赖必须仔细研究并计入系统不确定度。

\subsection{准 PDF 与 PDF 的一圈匹配}\label{sec:one_loop_matching}

要得到准 PDF $\widetilde{q}_R(x,P_z, p^R_z,\mu_R)$ 与光锥 PDF $q(x,\mu)$ 之间的匹配系数，可以在微扰论中计算非壳夸克矩阵元素。下文中，我们在朗道规范下分别对最小投影与 $\slashed p$ 投影进行计算。一般协变规范、一般洛伦兹结构的结果见附录~\ref{app:one-loop}。

最低阶的夸克准 PDF 为
\begin{equation}
\widetilde{q}^{(0)}(x)=\delta(1-x)\,.
\end{equation}
一圈阶则为
\begin{align}\label{eq:one_loop_quasiPDF}
\widetilde{q}^{(1)}(x,p,\rho)=\mbox{Tr}\left[\left(\left[\widetilde{f}_t(x,\rho)\right]_+\gamma^t+\left[\widetilde{f}_z(x,\rho)\right]_+\frac{p_t}{p_z}\gamma^z+\left[\widetilde{f}_p(x,\rho)\right]_+\frac{p_t\slashed{p}}{p^2}\right){\cal P}\right]\,.
\end{align}
其中 $\widetilde{f}_i$ 为
\begin{align}
\widetilde{f}_t(x,\rho)=\frac{\alpha_s C_F}{2\pi}\left\{
\begin{array}{lc}
\frac{8x^2(1-x)-x\rho(13-10x)+3\rho^2}{2(1-x)(1-\rho)(4x-4x^2-\rho)}+\frac{4x(2-x)-x\rho-3\rho}{4(1-x)(1-\rho)^{3/2}}\ln\frac{2x-1+\sqrt{1-\rho}}{2x-1-\sqrt{1-\rho}} & x>1\\
%
\frac{x(-7+4x)+3\rho}{2(1-x)(1-\rho)}+\frac{4x(2-x)-\rho(3+x)}{4(1-x)(1-\rho)^{3/2}}\ln\frac{1+\sqrt{1-\rho}}{1-\sqrt{1-\rho}} & 0<x<1\\
%
-\frac{8x^2(1-x)-x\rho(13-10x)+3\rho^2}{2(1-x)(1-\rho)(4x-4x^2-\rho)}-\frac{4x(2-x)-x\rho-3\rho}{4(1-x)(1-\rho)^{3/2}}\ln\frac{2x-1+\sqrt{1-\rho}}{2x-1-\sqrt{1-\rho}} & x<0
\end{array}\right.,
\end{align}
\begin{align}
\widetilde{f}_z(x,\rho)=\frac{\alpha_s C_F}{2\pi}\left\{
\begin{array}{lc}
\begin{array}{l}
\frac{-32x^2(1-x)^2(2x-1)-4x\rho(8-43x+65x^2-38x^3+8x^4)+\rho^2(5-41x+42x^2-8x^3)+2\rho^3(2-x)}{2(1-x)(1-\rho)^2(4x-4x^2-\rho)^2}\\
\quad+\frac{4-8x+8x^2+\rho(3-13x+4x^2)+2\rho^2}{4(1-x)(1-\rho)^{5/2}}\ln\frac{2x-1+\sqrt{1-\rho}}{2x-1-\sqrt{1-\rho}}
\end{array} & x>1\\
\frac{-5+15x-12x^2-2\rho(2-3x)}{2(1-x)(1-\rho)^2}+\frac{4-8x+8x^2+\rho(3-13x+4x^2)+2\rho^2}{4(1-x)(1-\rho)^{5/2}}\ln\frac{1+\sqrt{1-\rho}}{1-\sqrt{1-\rho}} & 0<x<1\\
\begin{array}{l}
-\frac{-32x^2(1-x)^2(2x-1)-4x\rho(8-43x+65x^2-38x^3+8x^4)+\rho^2(5-41x+42x^2-8x^3)+2\rho^3(2-x)}{2(1-x)(1-\rho)^2(4x-4x^2-\rho)^2}\\
\quad-\frac{4-8x+8x^2+\rho(3-13x+4x^2)+2\rho^2}{4(1-x)(1-\rho)^{5/2}}\ln\frac{2x-1+\sqrt{1-\rho}}{2x-1-\sqrt{1-\rho}}
\end{array} & x<0
\end{array}\right.,
\end{align}
\begin{align}
\widetilde{f}_p(x,\rho)=\frac{\alpha_s C_F}{2\pi}\left\{
\begin{array}{lc}
\frac{16x\rho(1-x)^2(1-6x)-2\rho^2(1-22x+26x^2-4x^3)-\rho^3(7-6x)}{2(1-\rho)^2(4x-4x^2-\rho)^2}+\frac{-\rho(8-12x+\rho)}{4(1-\rho)^{5/2}}\ln\frac{2x-1+\sqrt{1-\rho}}{2x-1-\sqrt{1-\rho}} & x>1\\
\frac{2-4x+\rho(7-8x)}{2(1-\rho)^2}+\frac{-\rho(8-12x+\rho)}{4(1-\rho)^{5/2}}\ln\frac{1+\sqrt{1-\rho}}{1-\sqrt{1-\rho}} & 0<x<1\\
-\frac{16x\rho(1-x)^2(1-6x)-2\rho^2(1-22x+26x^2-4x^3)-\rho^3(7-6x)}{2(1-\rho)^2(4x-4x^2-\rho)^2}-\frac{-\rho(8-12x+\rho)}{4(1-\rho)^{5/2}}\ln\frac{2x-1+\sqrt{1-\rho}}{2x-1-\sqrt{1-\rho}} & x<0
\end{array}\right.,
\end{align}
以及
\begin{align}
\rho=\frac{-p^2-i\epsilon}{p_z^2}
\end{align}
其中 $i\epsilon$ 给出了把 $\rho$ 从 $\rho<1$（闵氏）解析延拓到 $\rho>1$（欧氏）的规定。注意矢量流守恒保证顶点修正与波函数贡献可以合并成广义加（plus）函数~\cite{Stewart:2017tvs}。这些函数借助两个任意函数 $h(x)$ 与 $g(x)$ 定义：
\begin{align}
\int dx\left[{h(x)}\right]_+ g(x)&=\int dx\;h(x)\left[g(x)-g(1)\right]\,.
\end{align}

对于采用同样非壳红外调节的朗道规范下的光锥 PDF，树图贡献为
\begin{align}
q^{(0)}(x)=\delta(1-x)\,,
\end{align}
而 $\overline{\text{MS}}$ 方案中的一圈修正为
\begin{align}\label{eq:one_loop_lightconePDF}
&q^{(1)}(x,p,\mu)\nn\\
=&\mbox{Tr}\left[\left(\left[f_+\left(x,\frac{\mu^2}{p^2}\right)\right]_+\gamma^++\left[f_p(x)\frac{p^+\slashed{p}}{p^2}\right]_+\right) {\cal P}\right]\,.
\end{align}
其中
\begin{align}
f_+\left(x,\frac{\mu^2}{p^2}\right)=&\frac{\alpha_s C_F}{2\pi}\theta(x)\theta(1-x)\left[
\frac{-5+10x-6x^2}{2(1-x)}\right.\nn\\
&\left.+\frac{1+x^2}{1-x}\ln\frac{\mu^2}{-x(1-x)p^2}\right]\,, \\
f_p(x)=&\frac{\alpha_s C_F}{2\pi}(1-2x)\theta(x)\theta(1-x)\,.
\end{align}

要把准 PDF 匹配到光锥 PDF，需要对裸准 PDF 取在壳极限（$p^2\to 0$ 即 $\rho\to0$）与大动量极限（$p_t\to p_z$）：
\begin{equation}
\widetilde{q}^{(1)}_B(x,\rho)=\widetilde{q}^{(1)}(x,(p_t\to p_z,\vec{p}_\perp,p_z),\rho\to 0).
\end{equation}
可以看到式~(\ref{eq:one_loop_quasiPDF}) 中正比于 $\gamma^t$ 与 $\gamma^z$ 的两项在大动量极限下都趋于光锥算符，它们的组合捕获了正确的共线行为。因此最小投影下的裸准 PDF 定义为拾取式~(\ref{eq:one_loop_quasiPDF}) 中 $\gamma^t$ 与 $\gamma^z$ 的系数：
\begin{align}
\widetilde{q}_B^{(1)}(x,\rho)\Big|_{mp}=\left[\widetilde{f}_t(x,\rho)+\widetilde{f}_z(x,\rho)\right]_+\bigg|_{\rho\to 0}\,.
\end{align}
对于光锥 PDF，最小投影使用式~(\ref{eq:one_loop_lightconePDF}) 中 $\gamma^+$ 的系数：
\begin{align}
q^{(1)}(x,p,\mu)\Big|_{mp}= f_+\left(x,\frac{\mu^2}{p^2}\right)_+\,.
\end{align}
于是裸匹配系数推导为
\begin{align} \label{eq:mpbarematching}
f_{1,mp}\left(x,\frac{p_z}{\mu}\right)_+= \widetilde{q}^{(1)}_B(x,\rho)\Big|_{mp}-q^{(1)}(x,p,\mu)\bigg|_{mp},
\end{align}
其中
\begin{align}
&f_{1,mp}\left(x,\frac{p_z}{\mu}\right)=\frac{\alpha_s C_F}{2\pi} \nn\\
&\times \left\{
\begin{array}{lc}
\displaystyle \frac{1+x^2}{1-x}\ln\frac{x}{x-1}+1 & x>1\\
\displaystyle \frac{1+x^2}{1-x}\ln\frac{4x(1-x)p_z^2}{\mu^2}-\frac{x(1+x)}{1-x} & 0<x<1\\
\displaystyle -\frac{1+x^2}{1-x}\ln\frac{x}{x-1}-1 & x<0
\end{array} \right. .
\end{align}

在 RI/MOM 方案中，准 PDF 用一个额外的反项进行重正化。我们发现在 $|x|\to\infty$ 极限下，只有 $\widetilde{f}_t(x,\rho)$ 以 $1/|x|$ 行为衰减。对 $x$ 积分后，这一项重现定域极限 $z=0$ 处的紫外发散。因此，自然的选择是拾取式~(\ref{eq:one_loop_quasiPDF}) 中的 $\gamma^t$ 项作为反项：
\begin{align}
\widetilde{q}_{CT}^{(1)}\left(x,r,\frac{p_z}{p_z^R}\right)\bigg|_{mp}&=\left[\left|\frac{p_z}{p_z^R}\right|f_{2,mp}\left(1+\frac{p_z}{p_z^R}(x-1),r\right)\right]_+\,.
\end{align}
这里 $r=\mu_R^2/(p_z^R)^2$，
\begin{align}
f_{2,mp}(x,r)=\widetilde{f}_t(x,r)=\frac{\alpha_s C_F}{2\pi}\left\{
\begin{array}{lc}
\frac{-3r^2+13rx-8x^2-10rx^2+8x^3}{2(r-1)(x-1)(r-4x+4x^2)}+\frac{-3r+8x-rx-4x^2}{2(r-1)^{3/2}(x-1)}\tan^{-1}\frac{\sqrt{r-1}}{2x-1} & x>1\\
\frac{-3r+7x-4x^2}{2(r-1)(1-x)}+\frac{3r-8x+rx+4x^2}{2(r-1)^{3/2}(1-x)}\tan^{-1}\sqrt{r-1} & 0<x<1\\
-\frac{-3r^2+13rx-8x^2-10rx^2+8x^3}{2(r-1)(x-1)(r-4x+4x^2)}-\frac{-3r+8x-rx-4x^2}{2(r-1)^{3/2}(x-1)}\tan^{-1}\frac{\sqrt{r-1}}{2x-1} & x<0
\end{array} \right. .
\end{align}
最终，式~\eqref{eq:fact} 因子化公式中的匹配系数 $C$ 推导为
\begin{align}\label{eq:matching_coeff}
C\left(x,r,\frac{p_z}{\mu},\frac{p_z}{p_z^R}\right)&=\delta(1-x)+\left[\widetilde{q}^{(1)}_B(x,\rho)-q^{(1)}(x,p,\mu)-\widetilde{q}_{CT}^{(1)}\left(x,r,\frac{p_z}{p_z^R}\right)\right]\Bigg|_{mp}+{\cal O}(\alpha_s^2)\nonumber\\
&=\delta(1-x)+\left[f_{1,mp}\left(x,\frac{p_z}{\mu}\right)-\left|\frac{p_z}{p_z^R}\right|f_{2,mp}\left(1+\frac{p_z}{p_z^R}(x-1),r\right)\right]_+ +{\cal O}(\alpha_s^2)\,.
\end{align}
这里耦合常数 $\alpha_s(\mu)$ 采用标准 $\overline{\text{MS}}$ 方案。注意反夸克分布通过设定 $q(y)=-\bar{q}(-y)$ 映射到区域 $-1<y<0$。

对于 $\slashed{p}$ 投影，裸准 PDF 为
\begin{align}
\widetilde{q}_B^{(1)}(x,\rho)\Big|_{\slashed{p}}=\left[\widetilde{f}_t(x,\rho)+\widetilde{f}_z(x,\rho)+\widetilde{f}_p(x,\rho)\right]_+\,.
\end{align}
采用类似投影的光锥 PDF 为
\begin{align}
q^{(1)}(x,p,\mu)\Big|_{\slashed{p}}=\left[f_+\left(x,\frac{\mu^2}{p^2}\right)+f_p(x)\right]_+\,.
\end{align}
在该投影下，裸准 PDF 的匹配系数与式~(\ref{eq:mpbarematching}) 一致：
\begin{align}
f_{1,\slashed p}\left(x,\frac{p_z}{\mu}\right)_+&= \widetilde{q}^{(1)}_B(x,\rho)\Big|_{\slashed p}-q^{(1)}(x,p,\mu)\bigg|_{\slashed p} \nn\\
&=f_{1,mp}\left(x,\frac{p_z}{\mu}\right)_+\,.
\end{align}
反项可通过取 ${\cal P}=\slashed{p}/(4p^t)$ 计算得到：
\begin{align}
f_{2,\slashed{p}}(x,r)=\frac{\alpha_s C_F}{2\pi}\left\{
\begin{array}{lc}
\frac{3-3r-2x}{2(r-1)(x-1)}+\frac{4rx-8x^2+8x^3}{(r-4x+4x^2)^2}+\frac{2-2r-rx+2x^2}{(r-1)^{3/2}(x-1)}\tan^{-1}\frac{\sqrt{r-1}}{2x-1} & x>1\\
\frac{3-3r-2x+4x^2}{2(r-1)(1-x)}+\frac{-2+2r+rx-2x^2}{(r-1)^{3/2}(1-x)}\tan^{-1}\sqrt{r-1} & 0<x<1\\
-\frac{3-3r-2x}{2(r-1)(x-1)}-\frac{4rx-8x^2+8x^3}{(r-4x+4x^2)^2}-\frac{2-2r-rx+2x^2}{(r-1)^{3/2}(x-1)}\tan^{-1}\frac{\sqrt{r-1}}{2x-1} & x<0
\end{array} \right.\,.
\end{align}
相应的 RI/MOM 匹配系数由把式~(\ref{eq:matching_coeff}) 中的“$mp$”替换为“$\slashed p$”得到，且 $f_{2,\slashed{p}}$ 与 $f_{2,mp}$ 之差在 $p_z^R=0$ 极限下消失。$\Gamma=\gamma^z$ 情形的匹配系数也已在附录 \ref{app:one-loop_gamma_z} 中给出。
